%======================================================================
% supplementary.tex - LG-PM Supplementary Information
% Standalone compile: pdflatex supplementary x2
% Submitted alongside: "Lesion-grounded prototype meta-learning for
% few-shot thyroid ultrasound: segmentation, diagnosis, and faithful
% explanations under domain shift" -- Journal of Intelligent Medicine
% and Healthcare (JIMH), Tech Science Press.
%======================================================================
\documentclass[11pt]{article}
\usepackage[a4paper,margin=2.4cm]{geometry}
\usepackage{amsmath,amssymb}
\usepackage{booktabs}
\usepackage{graphicx}
\usepackage{algorithm}
\usepackage{algpseudocode}
\usepackage{float}
\usepackage{xcolor}
\renewcommand{\thesection}{S\arabic{section}}
\renewcommand{\thetable}{S\arabic{table}}
\renewcommand{\thefigure}{S\arabic{figure}}
\renewcommand{\theequation}{S\arabic{equation}}
\renewcommand{\thealgorithm}{S\arabic{algorithm}}
\newcommand{\lgShrink}[1]{\resizebox{\textwidth}{!}{#1}}

\title{\textbf{Supplementary Material}\\[4pt]
\normalsize for: Lesion-grounded prototype meta-learning for few-shot thyroid
ultrasound: segmentation, diagnosis, and faithful explanations under
domain shift\\[4pt]
\small S.~M.~Nihal Ahmed, Afrim Hossen Khan and Sabikun Nahar Sinthia\\[2pt]
\small Department of Software Engineering, Daffodil International University, Dhaka, Bangladesh\\[4pt]
\small \emph{Journal of Intelligent Medicine and Healthcare}}
\author{}
\date{}

\begin{document}
\maketitle

\noindent\emph{This document preserves, verbatim, the methodological
detail moved out of the 16-page main text. Each block is labelled and
cross-referenced to the title of the main-text subsection that cites
it. Main-text equations and section numbers referenced here are those
of the compressed manuscript.}

%======================================================================
% sm_methods.tex - Supplementary Methods (LG-PM 16-page split, 2026-08-23)
% Blocks below are VERBATIM extractions from proposed_methodology.tex
% (SM-CANDIDATE flags #1-#13, operator commission 2026-08-23).
% Main-text pointer labels use subsection titles, not numbers.
%======================================================================

\section{Supplementary data, annotation, and de-duplication detail}\label{sm:data}

\noindent\emph{S1.1 Per-image recorded fields (referenced from main text, Section `Data sources and annotation roles').}

The data pipeline records the following fields for every image:

\begin{itemize}
  \item dataset and domain identifier;
  \item patient, nodule, or examination identifier, where available
  (TN3K, TNUI-2021, and DDTI do not provide usable patient linkage in the
  released files - a recorded limitation);
  \item benign/malignant label and its source;
  \item mask, bounding box, or no localisation annotation;
  \item image dimensions, colour mode, and acquisition metadata, where
  available;
  \item split assignment and episode membership; and
  \item a cryptographic (SHA-256) image hash for duplicate checking.
\end{itemize}

\noindent\emph{S1.2 Annotation hierarchy (referenced from main text, Section `Data sources and annotation roles').}

The following annotation hierarchy is maintained:

\begin{enumerate}
  \item pixel masks supervise or evaluate segmentation;
  \item bounding boxes provide weaker localisation annotation and define the
  oracle region of interest where masks are absent (TN5000);
  \item image- or case-level pathology labels supervise malignancy
  classification; and
  \item TI-RADS or radiologist-derived labels are reported separately from
  pathology-confirmed outcomes.
\end{enumerate}

\noindent\emph{S1.3 Label parsers: full precedence specification (referenced from main text, Section `Data sources and annotation roles'; the TN5000 malignant-if-any rule and the TNUI polarity assumption-with-evidence are retained in the main text).}

\paragraph{Label provenance and harmonisation.}
Classification labels enter the pipeline through three documented parsers:

\begin{description}
  \item[TN5000 (source).] Pascal-VOC XML with object class names
  \texttt{"0"} (benign) and \texttt{"1"} (malignant). The image label is
  \emph{malignant-if-any} across annotated objects, and the oracle ROI box
  is the union of the boxes of the objects that decide the label. The
  official \texttt{ImageSets/Main} split files are used as released (the
  combined \texttt{trainval} file is not used).
  \item[TNUI-2021 (external).] Per-split CSV files (\texttt{ID},
  \texttt{Labels}; numeric 0/1) with identical schemas across splits. The
  polarity ($1 = $ malignant) is not stated in public documentation. We
  adopt this polarity as an \emph{assumption-with-evidence}: the
  alternative polarity
  would force in-domain-validated features to anti-classify on every TNUI
  cohort, and cohort sizes and prevalence are consistent with the
  published dataset description. The assumption is recorded in the
  limitations rather than silently absorbed.
  \item[DDTI (external).] Per-case XML diagnoses are parsed by an
  element-aware parser with a fixed precedence: literal
  \emph{benign}/\emph{malignant} strings override numeric diagnosis codes;
  diagnosis codes map $1 \rightarrow$ benign and $3 \rightarrow$
  malignant (all other codes unknown); TI-RADS annotations map
  $1$--$3 \rightarrow$ benign and $4^{(\ast)}$/$5 \rightarrow$ malignant.
  Evidence is aggregated malignant-if-any, mirroring the TN5000 rule, and
  cases without a parseable diagnosis are excluded and counted. A vote
  profile (how many cases decided by each evidence class) is printed in
  every run for audit.
\end{description}

\noindent\emph{S1.4 Exact-duplicate policy mechanics (referenced from main text, Section `Data splitting, de-duplication, and external-validation protocol').}

\paragraph{Exact-duplicate policy.}
Every image is hashed (SHA-256). Within each dataset, exact duplicate files
are removed once per dataset, keeping the occurrence in the highest-priority
partition (TN3K: trainval $>$ test; TNUI-2021: train $>$ val $>$ test).
Dropped duplicates are archived in a separate manifest.
Cross-dataset hashing runs against every pool that donates support cases (TN3K
training pool) and between every pair of classification cohorts. Exclusions
in later phases replicate the exclusion policy locked by earlier phases,
with a recompute fallback. Each experiment notebook independently
re-derives the clean
manifests from the raw dataset files. When the audit manifests from the
dataset-audit phase are mounted, the re-derived manifests are verified
against them at hash level. For classification phases, the TN5000 manifest must additionally
re-derive the exact locked split sizes before any model is trained or
applied.


\section{Supplementary model and training configuration}\label{sm:config}

\noindent\emph{S2.1 Stage-1 supervised reference: optimisation and inference (referenced from main text, `Stage 1 optimisation').}

\subsubsection{Stage 1 optimisation and inference details (supervised
 reference)}


The Stage 1 reference is a ResNet34-encoder U-Net. We train it with AdamW
(learning rate $10^{-3}$, weight decay $10^{-4}$), mixed precision, and
ReduceLROnPlateau scheduling (factor 0.5, patience 4), for at most 60 epochs
with early stopping (patience 12). An exponential moving average (EMA) of
the weights (decay 0.999) is maintained; each epoch selects the better of
raw and EMA weights by validation Dice. Training-time augmentation
(albumentations) comprises horizontal flip, affine scaling/translation
(0.85--1.15, $\pm$5\%) with rotation ($\pm10^{\circ}$),
brightness/contrast jitter ($\pm$0.25), Gaussian (speckle-like) noise,
mild Gaussian blur, and elastic deformation. At inference, we select a sigmoid
operating threshold on the internal validation split only
(sweep $0.10$--$0.70$, step $0.02$) and freeze it before any test or external
evaluation. Test-time augmentation (TTA; identity, horizontal flip,
$\pm5^{\circ}$ rotation; probability-averaged) is reported alongside
single-view inference. The tuned threshold value is reported in Results.


\noindent\emph{S2.2 Stage-2 episodic meta-training regimen (referenced from main text, `Stage 2: Episodic decoder meta-training').}

\subsubsection{Stage 2: Episodic decoder meta-training (single source
 domain)}

The encoder is frozen. Each training episode samples nested supports
($K_m$ re-drawn uniformly from $\{1,5,10,20\}$) and eight queries from the
source training pool. We apply light \emph{joint} augmentation (consistent
flip, $\pm10^{\circ}$ rotation, and per-image gain across all episode
images and masks), compute multi-scale prototypes with the frozen encoder,
and minimise the query loss with respect to the decoder parameters only.
Optimisation uses AdamW ($10^{-3}$, 300 episodes per epoch) with early
stopping on mean query Dice over 100 validation episodes (patience 8). Joint episodic meta-training
across multiple source domains is a designated extension; the primary
protocol trains on TN3K and characterises transfer to TNUI-2021 and DDTI
under the zero-shot and mask-shot conditions defined above.


\noindent\emph{S2.3 Linear-head training configuration (referenced from main text, `Classification heads').}

\paragraph{Linear heads (primary).}
Each regime's descriptor is standardised with the training-split mean and
standard deviation. We train a logistic linear head (a single affine layer)
with AdamW (learning rate $3\times10^{-3}$, weight decay
$10^{-2}$), binary cross-entropy with positive-class weighting
$(n_{\mathrm{neg}}/n_{\mathrm{pos}}$ on the training split), mini-batches
of 512, at most 300 epochs, and early stopping on validation AUROC
(patience 25, improvement margin $10^{-4}$, seed 42). The only quantity
that varies across regimes is \emph{what is pooled}; standardisation,
optimiser, and selection protocol are identical, so R0--R1p contrasts
isolate grounding and R1p--R1g contrasts isolate segmentation error.


\noindent\emph{S2.4 Whole-image ERM baseline configuration (referenced from main text, `Classification heads').}

\paragraph{Whole-image ERM baseline (the fair comparator).}
The comparator is whole-image empirical risk minimisation (ERM). An
ImageNet-pretrained ResNet34 with a single-logit head is fine-tuned
end-to-end on TN5000 images (BCE with positive-class weighting; AdamW,
learning rate $10^{-4}$, weight decay $10^{-4}$; cosine annealing over
$15$ epochs; early stopping on validation AUROC, patience 4; light
augmentation: horizontal flip, $\pm7^{\circ}$ rotation, gain jitter
$0.8$--$1.2$; seed 42). This row must be strong --- it is the conventional
alternative a reviewer compares the grounded pipeline against, both
in-domain and externally. An uncertainty-aware (shrunk-covariance)
prototype scorer is designated as an extension and is not part of the
reported heads.


\noindent\emph{S2.5 Calibration configuration (referenced from main text, `Stage 4: Calibration, operating points, and model selection').}

\subsubsection{Stage 4: Calibration, operating points, and model selection}

The source validation split (TN5000 val for classification; TN3K internal
val for segmentation) is used for: early stopping; the segmentation
operating-threshold sweep (validation only, then frozen); TTA
configuration; decoder checkpoint selection; and, for classification,
temperature scaling and the decision threshold. We fit a single scalar temperature
per model by grid search over log-spaced candidates
$T\in[0.05, 20]$, minimising validation negative log-likelihood on the
model's validation logits. The operating threshold per model is the
validation point of maximal sensitivity subject to specificity
$\geq 0.90$ on \emph{raw} validation probabilities. Both are frozen before
any test or external evaluation. No target-domain annotation is used for
any of these decisions.


\noindent\emph{S2.6 Feature-bank per-case quality fields (verbatim from the feature-bank governing convention; referenced from main text, `Lesion--context descriptor regimes').}

Cached per case alongside the five descriptors: empty-mask flag, predicted area fraction, mean foreground probability, a $96\times96$ mask thumbnail.

\section{Extended formulation}\label{sm:formulation}

\noindent\emph{S3.1 Episode and encoder definitions (referenced from main text, `Episode formulation' and `Shared multi-scale encoder').}

A segmentation episode pairs a support set with a query set:

\begin{equation}
  \mathcal{T}_{\mathrm{seg}}^{(d)}
  = \left(\mathcal{S}_{\mathrm{seg}}^{(d)},
  \mathcal{Q}_{\mathrm{seg}}^{(d)}\right),
  \label{eq:seg_task}
\end{equation}

The support set contains $K_m$ image--mask pairs and the query set $Q_m=8$ disjoint images:

\begin{equation}
  \mathcal{S}_{\mathrm{seg}}^{(d)}
  = \left\{(x_i,m_i)\right\}_{i=1}^{K_m},
  \qquad
  \mathcal{Q}_{\mathrm{seg}}^{(d)}
  = \left\{(x_j,m_j)\right\}_{j=1}^{Q_m},
  \qquad m_i\in\{0,1\}^{H\times W}.
  \label{eq:seg_support_query}
\end{equation}

The shared encoder maps an input image to multi-scale feature maps:

\begin{equation}
  \left\{F^l(x)\right\}_{l=1}^{L}=E_\theta(x),
  \qquad F^l(x)\in\mathbb{R}^{C_l\times H_l\times W_l}.
  \label{eq:multiscale_encoder}
\end{equation}

\noindent\emph{S3.2 Uncertainty-aware segmentation prototypes (pre-specified extension, not executed; equations referenced from main text, `Support-conditioned segmentation prototypes').}

\subsubsection{Uncertainty-aware segmentation prototypes (extension)}

A point prototype can be unstable when $K_m=1$ or when the support images are
not representative of the query domain. As a designated extension, a
channel-wise support variance may be estimated:

\begin{equation}
  v_c^l
  = \frac{1}{K_m}\sum_{i=1}^{K_m}
    \left(p_{c,i}^{l}-p_c^l\right)^2,
  \qquad c\in\{f,b\},
  \label{eq:prototype_variance}
\end{equation}

where $p_{c,i}^l$ is the prototype from support image $i$, stabilised by a
source-validation shrinkage coefficient $\rho$:

\begin{equation}
  \tilde v_c^l
  =(1-\rho)v_c^l+\rho\,\mathrm{mean}(v_c^l),
  \qquad 0\leq\rho\leq1,
  \label{eq:variance_shrinkage}
\end{equation}

with the uncertainty-adjusted similarity

\begin{equation}
  \tilde s_c^l(x_q,u)
  =-
  \frac{\left(F^l(x_q,u)-p_c^l\right)^2}
       {\tilde v_c^l+\epsilon},
  \label{eq:uncertain_similarity}
\end{equation}

with channel-wise division and aggregation across channels. This variant is
not part of the executed pipeline: the point-prototype model is the primary
and only reported configuration. We retain the variance-weighted form here
only as a pre-specified extension.


\noindent\emph{S3.3 Dense segmentation metric definitions (referenced from main text, `Evaluation metrics').}

Dice and intersection-over-union (IoU) are computed per case as:

\begin{equation}
  \mathrm{Dice}
  =\frac{2|M\cap\hat M|}{|M|+|\hat M|},
  \qquad
  \mathrm{IoU}
  =\frac{|M\cap\hat M|}{|M\cup\hat M|}.
  \label{eq:dice_iou_metric}
\end{equation}

\section{Pre-registered protocols: full specifications}\label{sm:protocols}

\noindent\emph{S4.1 Phase-4 joint fine-tuning: full regimen (referenced from main text, `Pre-registered joint partial-supervision phase'; the v1 gate failure and the 0.8066 guard are retained in the main text).}

Training follows protocol v2 (amended 2026-08-21, before any family-$G$/$H$
evaluation). \emph{Phase W} warms the classification head on the
\emph{frozen} trunk (AdamW $10^{-4}$ / $10^{-4}$, cosine annealing, 25
epochs, classification batches only), placing the head near its conditional
minimum before the trunk is touched. \emph{Phase J} then unfreezes the trunk
and jointly fine-tunes on the interleaved batches at $10^{-5}$, an order of
magnitude below the ERM regimen, with cosine annealing, at most 15 epochs,
patience 4, seed 42. Checkpoint selection maximises source validation AUROC
\emph{subject to a segmentation guard}: TN3K-validation Dice
$\geq 0.8066$ (the pre-registered non-collapse gate constant plus a $0.01$
margin); the test split never participates in selection. The guard exists
because protocol v1 failed this gate. Under a single-stage joint fine-tune
at $10^{-4}$ with AUROC-only
selection, the classification gradient of the freshly initialised head rotated the
shared encoder from the first epoch: saved-model TN3K-test Dice was 0.6686
against a gate of 0.7966. The harness itself was validated by a
pristine-checkpoint
control that reproduced the locked reference exactly. Should no joint epoch
satisfy the guard, the best-Dice state is retained for diagnosis and the
non-collapse gate is allowed to fail. The joint phase is then reported as a
negative result rather than analysed. Both arms share identical
batch/augmentation streams (\texttt{set\_seed(42)} first, warmup consuming a
fixed, data-independent number of draws per arm) and differ only in the
classification pooling region. Evaluation
re-uses the locked caches wherever the protocol permits. The in-domain
family $G$ (pre-declared, Bonferroni within family) comprises: G1
$\mathrm{JOINTg}-\mathrm{ERM}$ AUROC (is the in-domain fine-tuning gap
closed?), G2 $\mathrm{JOINTg}-\mathrm{R1p}$ (frozen versus joint
grounding), G3 $\mathrm{JOINTg}-\mathrm{JOINT\textrm{-}img}$ (grounding
versus multi-tasking), G4 sensitivity@spec90 $\mathrm{JOINTg}-\mathrm{ERM}$
(raw-raw). The joint model then enters the frozen external battery defined by the main-text transfer protocol,
%% SM-EDIT: cross-document pointer (main-text target). with thresholds and
temperature fitted on its own TN5000 validation probabilities only.
The TNUI-family operating-point behaviour (pre-specified as primary) and
the DDTI descriptive rows are reported exactly as for earlier models.
Acceptance gates fixed here: the TN5000 manifest must re-derive the locked
split sizes; the joint segmentation branch must not collapse relative to
the locked Stage-1 references; and the ERM sanity gate against the locked
ERM probabilities applies unchanged. A pre-registered uncertainty
component derives predictive entropy and probability margins from the
calibrated probabilities and reports coverage-versus-malignant-FNR
selective-prediction curves on the source test and the pre-declared pool.


\noindent\emph{S4.2 Phase-5 audit gate: full A1--A5 mechanics AND the verbatim pre-registered family J definition (referenced from main text, `Case-level external validation on ThyUS2Path'; gate thresholds, the WITHHELD rule, the amendment date, and the J1 primary contrast are retained in the main text).}

\paragraph{Audit-gated inference (A1--A5).} No metric cell may run before
the audit gate, and no inferential statement is unsealed before the
licence. The checks execute in order:
\begin{description}
  \item[A1 label coverage.] A two-table per-batch prefix join maps each
  image identifier (first underscore-delimited token, namespaced by
  batch to forestall cross-batch collisions) onto its batch's label
  table (\texttt{histo\_label} $\in \{0, 1\}$). Per-batch coverage must
  reach 95\%, and the surviving join is cross-checked against the
  corpus's own image map.
  \item[A2 case assembly.] A case is the set of images sharing its
  batch-namespaced patient key. The primary analytic unit is the case,
  scored by the unweighted mean of its per-image calibrated
  probabilities; a first-listed-image view is pre-declared as a
  descriptive sensitivity analysis.
  \item[A3 overlap screening.] SHA-256 exact hashing plus perceptual-hash
  near-duplicate detection against every support or training pool
  (TN5000, TNUI-2021, the TN3K training pool); a single matching image
  excludes the \emph{whole} case (conservative policy).
  \item[A4 cohort gates.] Minimum assembled size (300 cases) and a
  malignancy-prevalence window of $[0.20, 0.85]$.
  \item[A5 polarity attestation.] The operator inspects a seeded grid of
  benign- and malignant-labelled images against the corpus documentation
  and attests that label polarity matches ultrasonic appearance; only
  then is the inference licence set.
\end{description}
The pre-registered family is \emph{computed} under the sealed discipline
but reported as \textsc{withheld} (descriptive statistics only) unless
A1--A4 pass computationally \emph{and} A5 is attested. The licence
upgrade flips the significance column only and recomputes nothing. The
amendment fixing the two-table join, the case-mean primary unit, and the
conservative whole-case exclusion policy was adopted 2026-08-22, strictly
before any target-side forward pass.

\paragraph{Pre-registered family J} (case level, 5{,}000 paired bootstrap
resamples, two-sided empirical $p$, Bonferroni
$\alpha = 0.05/4 = 0.0125$; raw-raw convention; thresholds and
temperatures from TN5000 validation only, nothing refitted):
\textbf{J1} sensitivity@spec90 R1p $-$ ERM, the primary external
operating-point replication, with McNemar discordants on malignant
cases; \textbf{J2} AUROC R1p $-$ ERM (rank transfer, secondary);
\textbf{J3} sensitivity@spec90 JOINT-g $-$ R1p (external analogue of
H2); \textbf{J4} AUROC R1p $-$ R0 (grounding versus whole image).
Oracle regimes (R1g/R2g) are excluded by pre-declaration because
ThyUS2Path provides no lesion annotation.


\noindent\emph{S4.3 Phase-6 explainability battery: full axis specifications, perturbation recipes, cohort paragraph, and the verbatim pre-registered family K definition (referenced from main text, `Pre-registered lesion-grounded explainability battery'; axis summaries, the ELA reproduction assertion, the cascade self-gate constant, cohort sizes, and the K1 primary contrast are retained in the main text).}

\paragraph{Axes.} (a) \emph{Plausibility}: explanation concentration in
the TN5000 expert box --- the \emph{localisation ratio} = in-box
explanation share divided by box-area share (segmentation-region
precision for descriptor heads; CAM energy for ERM), plus descriptive
pointing-game hits. The ThyUS2Path side is descriptive only (predicted-mask
agreement; no GT). (b) \emph{Faithfulness}: deletion (primary) and
insertion curves under a per-image Gaussian-blur baseline
($\sigma = 8$, 8 steps), raw probabilities (raw-raw convention), masks
held fixed at the original prediction (isolated descriptor sensitivity),
and one seeded random pixel ordering per image as control. The
\emph{faithfulness gap} is the trapezoidal area difference between the
random-ordering and attribution-guided curves. (c) \emph{Stability}:
seeded Gaussian noise ($\sigma = 0.02$) and affine jitter ($\pm 3^{\circ}$,
$\pm 5$\,px) with the full pipeline (masks) recomputed; Spearman
correlation of the clean and perturbed CAMs (affine maps inverse-warped
before comparison). (d) \emph{Randomisation sanity}: cascading
reinitialisation --- R1p head weights; ERM fc $\to$ fc+layer4 $\to$
fc+layer4+layer3 --- on a seeded 160-image probe; the battery self-gates:
mean $\lvert\rho\rvert < 0.2$ at the final cascade stage, asserted.

\paragraph{Cohorts.} TN5000-test (canonical Phase-3 chain, $n{=}950$, 693
malignant; expert boxes as GT) and a seeded ThyUS2Path case-first-image
panel ($n{=}600$, drawn from the licensed Phase-5 839-case manifest).
Metrics on this panel are label-free, so the battery is independent of
the Phase-5 licence.

\paragraph{Pre-registered family K} (image level; 5{,}000 paired bootstrap
resamples; two-sided empirical $p$; Bonferroni $\alpha = 0.05/4 = 0.0125$;
seed 42): \textbf{K1} GT-localisation ratio R1p $-$ ERM on TN5000-test
(PRIMARY plausibility); \textbf{K2} deletion faithfulness gap R1p $-$ ERM
(TN5000-test); \textbf{K3} deletion faithfulness gap R1p $-$ ERM
(ThyUS2Path panel; label-free); \textbf{K4} deletion faithfulness gap
JOINT-g $-$ JOINT-img (TN5000-test; explanation-space mirror of G3).
Everything else --- insertion curves, pointing games, CAM-variant
robustness, stability, cascade curves --- is descriptive with
pre-declared expectations.


\noindent\emph{S4.4 External classification transfer protocol: full gate specifications, tolerances, replication gate, and pooling/power policy (referenced from main text, `External classification transfer protocol'; the two-sentence gate-logic summary is retained in the main text).}

\paragraph{Frozen objects.}
Encoder, fusion decoder, fixed support episode (identifier-asserted against
the archive), descriptor code, head weights, operating thresholds, and
temperatures are byte-identical to the source phase.

\paragraph{Tolerance gates (checked before any external table).}
\begin{enumerate}
  \item \textbf{Support-episode gate:} episode identifiers must equal the
  archived identifiers (seed-drift assert).
  \item \textbf{Head-provenance gate:} heads are loaded from the locked
  phase caches. Where a head must be retrained from cached features under
  the identical protocol (the random-number generator, RNG, is not exactly
  reproducible), its source validation AUROC must lie within $\pm 0.01$ of
  the locked value; otherwise the run stops.
  \item \textbf{ERM sanity gate:} the baseline's source-test AUROC,
  recomputed from the archived probabilities/checkpoint, must match the
  locked value within $\pm 0.002$.
  \item \textbf{Segmentation-QC gate:} external predicted-mask Dice is
  reported against the locked Phase-2 zero-shot references before any
  classification analysis.
\end{enumerate}

\paragraph{Replication gate (governs pooling).}
When a previously locked external cohort is recomputed in a later session,
the new table must match the locked table within absolute tolerance
($0.0015$ with cached heads; $0.01$ with tolerated retrains) on every row.
If any row drifts beyond tolerance, all analyses depending on pooling are
\emph{withheld} (flagged \texttt{POOLED\_OK=False}; the sections report
withholding rather than crashing or silently recomputing).

\paragraph{Pooling policy.}
Exactly one pool is pre-declared: TNUI-2021 test + validation, with its own
test family (P1--P5, per-family Bonferroni). Later cohorts (TNUI-2021
train, DDTI) are analysed as individual families (E, F) with their own
Bonferroni screens. Cross-cohort evidence matrices (AUROC and
operating-point sensitivity across all cohorts) are descriptive only and
carry no inferential claims beyond the pre-declared families.

\paragraph{Power policy.}
External cohorts with very few benign cases yield wide intervals; the
pre-declared remedies are direction-consistency reporting across cohorts
and the single pooled analysis. Underpowered individual tests are labelled
as such rather than narrated as effects.


\section{Reproducibility artefact manifest}\label{sm:repro}

\noindent\emph{Referenced from main text, `Reproducibility and implementation controls'.}

Released or documented:

\begin{itemize}
  \item dataset versions and access locations;
  \item clean manifests re-derived from raw files and hash-verified
  against the audit manifests;
  \item duplicate-drop lists;
  \item the TN5000 manifest with locked split sizes and per-split
  prevalence;
  \item preprocessing and augmentation configurations;
  \item episode manifests, including the archived fixed support episode;
  \item feature-bank caches with quality-control fields;
  \item cached model probabilities for every reported table;
  \item label parsers and the DDTI vote profile;
  \item the TNUI label-polarity record (assumption-with-evidence);
  \item locked result tables cited by later phases as tolerance
  references;
  \item the audit-hash manifest of every external phase;
  \item seeds and software versions;
  \item model checkpoints (the Stage 1 reference checkpoint and the
  fusion decoder are versioned with the repository artefacts); and
  \item scripts for metrics and statistics.
\end{itemize}

\section{Algorithmic summary}\label{sm:algorithm}

\noindent\emph{Algorithm S1, referenced from main text, `Algorithmic summary'.}

\begin{algorithm}[!ht]
\caption{Lesion-Grounded Few-Shot Segmentation--Classification Framework}
\label{alg:lgpm}
\begin{algorithmic}[1]
\Require Source segmentation domain with masks; pathology-labelled source
classification domain; external cohorts; budgets $K_m,K_c$; query counts;
seeds
\Ensure Segmentation masks; malignancy predictions; calibration and
operating points; explanations
\State Hash and de-duplicate images; rebuild and verify manifests; group by
patient/nodule where identifiers exist.
\State \textbf{Stage 1:} Train the supervised ResNet34-UNet reference on the
source training pool (Tversky objective, EMA, val-only threshold sweep,
frozen threshold).
\State \textbf{Stage 2:} Freeze the encoder; train the fusion decoder
episodically on source episodes with nested budgets and joint augmentation.
\State Compute multi-scale fg/bg prototypes
(the main-text multi-scale foreground/background prototype equations);
%% SM-EDIT: cross-document pointer (main-text target).
predict query masks via the main-text query-mask decoder equation or the direct variant
%% SM-EDIT: cross-document pointer (main-text target)..
\State Evaluate in-domain, zero-shot, and mask-shot segmentation over 200
episodes each, all budgets on matched queries; bootstrap episode means.
\State \textbf{Stage 3:} Archive the fixed support episode; segment all
classification cohorts zero-shot; build and cache the five descriptor
regimes; train linear heads, the prototype budget sweep, and the ERM
baseline on the source domain only.
\State \textbf{Stage 4:} Freeze per-model temperature and the
specificity-$0.90$ operating threshold from source validation (raw-raw
convention).
\State \textbf{Stage 5:} Apply the frozen pipeline to each external cohort;
enforce tolerance, sanity, and replication gates; run the pre-declared
paired families and the single pre-declared pool.
\State \textbf{Phase 4 (pre-registered):} joint partial-supervision
fine-tune; family $G$ in-domain; external battery under the same gates.
\State Evaluate lesion alignment, faithfulness, stability, and randomisation
of explanations on internal and external cases.
\end{algorithmic}
\end{algorithm}

%======================================================================
% sm_results.tex - Supplementary Results (LG-PM 16-page split).
% All tables VERBATIM from results_lgpm_draft.tex (whole tables), or
% verbatim ROWS relocated from the pre-registered family tables
% (assembler: build_sm_results.py; commission 2026-08-23, stage 2).
% Inclusion order FIXES the SM table numbers S1..S22 cited from the
% compressed main text (sec16_results.tex).
%======================================================================

\section{Supplementary results}\label{sm:results}

\noindent\emph{This section preserves, verbatim, the per-cohort and
operational tables displaced from the 16-page main text, and -- in
Tables S5 and S22 -- the secondary and tertiary pre-registered
contrasts relocated row-for-row from the main-text family tables.
Each block notes the main-text subsection that cites it.}

\noindent\emph{Table S1, referenced from main text, `Cohorts, audits, and frozen operating conventions'.}

\begin{table}[!ht]
\centering\small
\caption{Frozen-pipeline segmentation QC in classification phases (locked,
Phase-1/2 pipeline at the fixed $K_m{=}20$ support episode; Dice vs cohort
masks at the locked 0.5 binarisation). ---: not available --- predicted-mask
area was not recorded by the Phase-3d/3e banks (Dice/empty only), and the
TNUI-train empty-rate line was not archived; descriptive QC only, both
recoverable on re-run.}
\label{tab:segqc}
\begin{tabular}{lccc}
\toprule
Cohort & mean Dice & empty-mask rate & mean mask area \\
\midrule
TN5000 (all splits) & n/a (expert boxes; no pixel masks) & 1.1\% val / 1.2\% test & 0.037--0.038 \\
TNUI-2021 test  & 0.7046 & 0.0\% & 0.033 \\
TNUI-2021 val   & 0.7307 & 1.5\% & --- \\
TNUI-2021 train & 0.7241 & --- & --- \\
DDTI (labelled subset) & 0.6132 & 3.7\% & --- \\
\bottomrule
\end{tabular}
\end{table}

\noindent\emph{Table S2, referenced from main text, `Supervised segmentation reference and external gap'.}

\begin{table}[!ht]
\centering\footnotesize
\caption{Phase-1 supervised segmentation reference (locked; from
\texttt{aggregate\_metrics.csv}). 95\% case-bootstrap CIs for Dice; HD95 is the
median over cases with both masks non-empty (n given; letterboxed
$384^2$ pixels). TNUI per-split TTA Dice: train 0.7894, val 0.8029.}
\label{tab:phase1}
\begin{tabular}{llccccc}
\toprule
Set & TTA & Dice (95\% CI) & IoU & Sens & Prec & HD95 med (n) \\
\midrule
TN3K test ($n{=}614$) & yes & 0.8466 (0.8336--0.8589) & 0.7603 & 0.9108 & 0.8240 & 11.16 (613) \\
                      & no  & 0.8441 (0.8310--0.8565) & 0.7572 & 0.9149 & 0.8165 & 11.31 (613) \\
TNUI-2021 all ($n{=}1371$) & yes & 0.7877 (0.7778--0.7972) & 0.6791 & 0.9117 & 0.7329 & 9.25 (1366) \\
                      & no  & 0.7766 (0.7663--0.7867) & 0.6665 & 0.9150 & 0.7160 & 10.00 (1368) \\
TNUI-2021 test ($n{=}138$) & yes & 0.7464 (0.7098--0.7805) & 0.6315 & 0.9048 & 0.6837 & 10.63 (138) \\
                      & no  & 0.7367 (0.6977--0.7713) & 0.6211 & 0.9072 & 0.6699 & 11.01 (138) \\
DDTI ($n{=}635$)      & yes & 0.7396 (0.7213--0.7561) & 0.6288 & 0.8324 & 0.7327 & 22.00 (629) \\
                      & no  & 0.7431 (0.7258--0.7589) & 0.6301 & 0.8471 & 0.7252 & 23.22 (629) \\
\bottomrule
\end{tabular}
\end{table}

\noindent\emph{Table S3, referenced from main text, `Few-shot support-conditioned segmentation'.}

\begin{table}[!ht]
\centering\footnotesize
\caption{Phase-2 direct (parameter-free) variant, query Dice (locked). IoU for
the decoder at $K{=}20$: 0.6873 / 0.5808 / 0.5471 / 0.5598 / 0.5602 across the
five conditions in table order.}
\label{tab:phase2_direct}
\resizebox{\textwidth}{!}{%
\begin{tabular}{lcccc}
\toprule
Condition & $K{=}1$ & $K{=}5$ & $K{=}10$ & $K{=}20$ \\
\midrule
TN3K test (in-domain)  & 0.6205 (0.5987--0.6432) & 0.6778 (0.6642--0.6911) & 0.6788 (0.6656--0.6922) & 0.6751 (0.6626--0.6877) \\
TNUI test (zero-shot)  & 0.4744 (0.4510--0.4975) & 0.4943 (0.4727--0.5161) & 0.4964 (0.4762--0.5170) & 0.4985 (0.4795--0.5174) \\
TNUI test ($K$-shot)   & 0.4906 (0.4657--0.5144) & 0.5965 (0.5823--0.6104) & 0.6175 (0.6040--0.6297) & 0.6277 (0.6156--0.6399) \\
DDTI (zero-shot)       & 0.5111 (0.4888--0.5332) & 0.5637 (0.5480--0.5787) & 0.5665 (0.5519--0.5800) & 0.5682 (0.5544--0.5813) \\
DDTI ($K$-shot target) & 0.5789 (0.5616--0.5965) & 0.6429 (0.6323--0.6545) & 0.6507 (0.6398--0.6624) & 0.6529 (0.6418--0.6642) \\
\bottomrule
\end{tabular}%
}
\end{table}

\noindent\emph{Table S4, referenced from main text, `In-domain classification with frozen descriptors'.}

\begin{table}[!ht]
\centering\scriptsize
\setlength{\tabcolsep}{3pt}
\caption{TN5000 test classification (locked; $n{=}950$, malignant 693).
ECE 15-bin; Brier and ECE-post use validation-fitted temperature;
sens@spec90 is evaluated at the frozen raw validation threshold (raw-raw
convention). CIs: 95\% case bootstrap (2,000 replicates).}
\label{tab:phase3}
\begin{tabular}{lcccccc}
\toprule
Model & AUROC (95\% CI) & AUPRC & sens@spec90 & Brier & ECE pre/post & $T_{\mathrm{val}}$ \\
\midrule
ERM whole-image (ResNet34 ft) & 0.9401 (0.9232--0.9553) & 0.9738 & 0.8023 & 0.0878 & 0.0893/0.0634 & 1.969 \\
R0 image | linear             & 0.8534 (0.8271--0.8798) & 0.9358 & 0.6479 & 0.1520 & 0.1278/0.1244 & 0.900 \\
R1p lesion | linear           & 0.8902 (0.8670--0.9122) & 0.9519 & 0.7273 & 0.1515 & 0.1788/0.1734 & 0.847 \\
R2p lesion+ring | linear      & 0.8829 (0.8584--0.9061) & 0.9464 & 0.6869 & 0.1538 & 0.1721/0.1705 & 0.956 \\
R1g GT box | linear (oracle)  & 0.9105 (0.8899--0.9299) & 0.9635 & 0.8254 & 0.1260 & 0.1339/0.1302 & 0.927 \\
R2g GT box+ring | linear (oracle) & 0.9029 (0.8807--0.9239) & 0.9569 & 0.7763 & 0.1449 & 0.1750/0.1721 & 0.927 \\
\bottomrule
\end{tabular}
\end{table}

\noindent\emph{Table S5: secondary and tertiary pre-registered contrasts
of families T, G, H, and J, relocated row-for-row from main-text Tables
`Phase-3b paired tests', `Phase-4 pre-registered families', and
`Phase-5 pre-registered family J' (referenced from the corresponding
main-text subsections). Cohort, resample count, and Bonferroni bar per
family follow each family header row; conventions as in the main-text
tables ($\Delta$ = A$-$B; empirical $p$ floored at $p{<}0.0004$).}

\begin{table}[!ht]
\centering\small
\caption{Supplementary secondary/tertiary pre-registered contrasts
(relocated rows; locked values). Family T: TN5000 test, 5,000 paired
bootstrap resamples, $\alpha{=}0.0083$. Family G: in-domain TN5000
test, 5,000 resamples, $\alpha{=}0.0125$. Family H: pooled TNUI
test+val ($n{=}413$, 392 positives), 2,000 resamples,
$\alpha{=}0.0167$. Family J: ThyUS2Path case level ($n{=}839$, 522
malignant), 5,000 resamples, $\alpha{=}0.0125$.}
\label{tab:sm_secondaries}
\resizebox{\textwidth}{!}{%
\begin{tabular}{clllccl}
\toprule
Test & Contrast & Metric & $\Delta$ (95\% CI) & $p$ & Sig. & Question \\
\midrule
\multicolumn{7}{l}{\textbf{Family T (secondaries), in-domain TN5000 test}}\\
T3 & R1g $-$ R1p & AUROC & +0.0202 (0.0056, 0.0349) & 0.0048 & YES & seg-error cost $\approx$2 pts \\
T4 & R1p $-$ R2p & AUROC & +0.0074 ($-$0.0023, 0.0176) & 0.1388 & no & ring adds nothing (H6 null) \\
T5 & R1p $-$ ERM & sens@spec90 & $-$0.0749 ($-$0.1126, $-$0.0390) & $<$0.0004 & YES & ERM wins op.\ point \\
T6 & R1g $-$ R1p & sens@spec90 & +0.0982 (0.0693, 0.1270) & $<$0.0004 & YES & oracle $>$ predicted ROI \\
\multicolumn{7}{l}{McNemar (T5): ERM-right/R1p-wrong 114 vs R1p-right/ERM-wrong 62, $\chi^2 = 14.78$.} \\
\midrule
\multicolumn{7}{l}{\textbf{Family G (G2--G4), in-domain}}\\
G2 & JOINT-g $-$ R1p & AUROC & $-$0.0205 ($-$0.0365, $-$0.0045) & 0.0124 & YES & joint vs frozen grounding --- frozen wins \\
G3 & JOINT-g $-$ JOINT-img & AUROC & +0.1489 (0.1136, 0.1848) & $<$0.0004 & YES & grounding beats multi-tasking \\
G4 & JOINT-g $-$ ERM & sens@spec90 & $-$0.2376 ($-$0.2785, $-$0.1962) & $<$0.0004 & YES & in-domain operating point --- ERM wins \\
\multicolumn{7}{l}{McNemar (G4, test positives): JOINT-g-right/ERM-wrong 41 vs ERM-right/JOINT-g-wrong 206,}\\
\multicolumn{7}{l}{$\chi^2 = 108.89$ ($p \ll 0.01$) --- $\approx$5:1 discordants against JOINT-g at the frozen point.}\\
\multicolumn{7}{l}{\textbf{Family H (H2--H3), pooled TNUI test+val (external)}}\\
H2 & JOINT-g $-$ R1p & sens@spec90 (pooled) & $-$0.0534 ($-$0.0969, $-$0.0102) & 0.011 & YES & frozen grounding transfers better \\
H3 & JOINT-g $-$ ERM & AUROC (pooled) & $-$0.0109 ($-$0.1442, 0.1322) & 0.858 & no & external rank (secondary) --- null \\
\midrule
\multicolumn{7}{l}{\textbf{Family J (J2--J4), ThyUS2Path case level (external)}}\\
J2 & R1p $-$ ERM      & AUROC       & $-$0.0401 ($-$0.0801, $-$0.0015) & 0.0436 & no  & rank transfer --- null at family bar \\
J3 & JOINT-g $-$ R1p  & sens@spec90 & $-$0.0978 ($-$0.1264, $-$0.0709) & $<$0.0004 & YES & frozen grounding transfers better (H2 analogue) \\
J4 & R1p $-$ R0       & AUROC       & $-$0.0322 ($-$0.0709, 0.0075)  & 0.1104 & no  & grounding vs whole image --- null \\
\bottomrule
\end{tabular}%
}
\end{table}

\noindent\emph{Table S6, referenced from main text, `In-domain classification with frozen descriptors'.}

\begin{table}[!ht]
\centering\small
\caption{Frozen operating points (locked): per-model thresholds chosen on raw
validation probabilities at specificity $\geq 0.90$, then applied unchanged to
the test split (raw-raw convention). The val$\to$test sensitivity transfer gap
is smallest for the oracle and grounded heads (R1g $+0.60$ pts, R1p $-1.71$
pts) and largest for ERM ($-5.33$ pts), foreshadowing the external
operating-point behaviour.}
\label{tab:spec90}
\begin{tabular}{lcccc}
\toprule
Model & threshold (frozen) & val spec & val sens & test sens@spec90 \\
\midrule
ERM            & 0.9233 & 0.9052 & 0.8556 & 0.8023 \\
R0 image       & 0.6708 & 0.9052 & 0.6417 & 0.6479 \\
R1p lesion     & 0.5323 & 0.9052 & 0.7444 & 0.7273 \\
R2p lesion+ring & 0.5789 & 0.9138 & 0.7278 & 0.6869 \\
R1g GT box (oracle) & 0.5194 & 0.9052 & 0.8194 & 0.8254 \\
R2g GT box+ring (oracle) & 0.4699 & 0.9052 & 0.8056 & 0.7763 \\
\bottomrule
\end{tabular}
\end{table}

\noindent\emph{Table S7, referenced from main text, `In-domain classification with frozen descriptors' (error-propagation quartiles).}

\begin{table}[!ht]
\centering\small
\caption{Error propagation (RQ4, locked; TN5000 test, $n{=}950$):
oracle-minus-predicted AUROC gap within predicted-mask-confidence quartiles
(237--238 cases each). gap = AUROC(GT-ROI) $-$ AUROC(pred-ROI); positive =
oracle better. The gap shows no monotone confidence relationship and is
descriptively largest at the highest confidence (single split; no
per-quartile CIs).}
\label{tab:propagation}
\begin{tabular}{lcccc}
\toprule
Confidence quartile & $n$ & AUROC pred-ROI & AUROC GT-ROI & gap \\
\midrule
Q1 (low)  & 238 & 0.8871 & 0.9081 & +0.0209 \\
Q2        & 237 & 0.8720 & 0.8834 & +0.0114 \\
Q3        & 237 & 0.8811 & 0.8769 & $-$0.0042 \\
Q4 (high) & 238 & 0.8433 & 0.9022 & +0.0589 \\
\bottomrule
\end{tabular}
\end{table}

\noindent\emph{Table S8, referenced from main text, `In-domain classification with frozen descriptors' (label-budget sweep).}

\begin{table}[!ht]
\centering\small
\caption{Class-prototype label-budget sweep (locked; TN5000 test).
Cosine-margin prototype classifier on frozen descriptors; AUROC mean $\pm$
SD (SD to 3 d.p.) over 25 independent per-class label draws at each budget
$K_c$; the \emph{all} row uses the full training split once (single draw).
Grounded heads dominate the image budget at every $K$ and saturate by
$K{=}50$; R0 does not.}
\label{tab:protosweep}
\begin{tabular}{lccccc}
\toprule
Budget $K_c$ & R0 image & R1p lesion & R2p lesion+ring & R1g GT box & R2g GT box+ring \\
\midrule
5   & 0.5058$\pm$0.050 & 0.5862$\pm$0.048 & 0.5778$\pm$0.044 & 0.6119$\pm$0.051 & 0.6066$\pm$0.033 \\
10  & 0.5307$\pm$0.045 & 0.6190$\pm$0.043 & 0.6058$\pm$0.036 & 0.6477$\pm$0.053 & 0.6290$\pm$0.030 \\
20  & 0.5398$\pm$0.044 & 0.6362$\pm$0.043 & 0.6246$\pm$0.037 & 0.6618$\pm$0.043 & 0.6484$\pm$0.032 \\
50  & 0.5586$\pm$0.034 & 0.6383$\pm$0.037 & 0.6335$\pm$0.033 & 0.6570$\pm$0.033 & 0.6478$\pm$0.024 \\
200 & 0.5827$\pm$0.032 & 0.6348$\pm$0.023 & 0.6294$\pm$0.023 & 0.6491$\pm$0.021 & 0.6441$\pm$0.015 \\
all & 0.5917 & 0.6379 & 0.6313 & 0.6508 & 0.6462 \\
\bottomrule
\end{tabular}
\par\smallskip
{\footnotesize Support-sampling variance shrinks with budget
($\sigma \approx 0.05$ at $K{=}5$ falling to $\approx 0.02$ at
$K{=}200$). Absolute levels sit far below the linear heads'
(0.59--0.66 vs 0.85--0.91), as expected for a cosine prototype
rule on frozen descriptors.}
\end{table}

\noindent\emph{Table S9, referenced from main text, `External classification battery' (TNUI-2021 test; per-cohort family Tc).}

\begin{table}[!ht]
\centering\scriptsize
\setlength{\tabcolsep}{3pt}
\caption{Phase-3c TNUI-test external classification (locked; $n{=}139$, 134
malignant). Thresholds frozen from TN5000-val raw probabilities (raw-raw);
Brier/ECE-post use source-val-fitted temperatures $T$ (frozen descriptors,
tolerance-gated retrained heads). AUPRC is near-uninformative at prevalence
0.964. External ECE is poor for \emph{every} model (0.24--0.43) ---
prevalence-shift miscalibration. Shift = in-domain minus external AUROC.}
\label{tab:phase3c}
\begin{tabular}{lccccccc}
\toprule
Model & AUROC (95\% CI) & AUPRC & sens@spec90 & Brier & ECE pre/post & $T_{\mathrm{src\,val}}$ & shift \\
\midrule
ERM            & 0.5172 (0.3296--0.6937) & 0.9745 & 0.5373 & 0.1963 & 0.2427/0.2674 & 1.969 & 0.423 \\
R0             & 0.5791 (0.3082--0.8154) & 0.9746 & 0.4104 & 0.3074 & 0.4318/0.4244 & 0.927 & 0.273 \\
R1p            & 0.7119 (0.5612--0.8768) & 0.9871 & 0.6940 & 0.2036 & 0.3173/0.3060 & 0.847 & 0.177 \\
R2p            & 0.6776 (0.3630--0.9058) & 0.9797 & 0.6493 & 0.2157 & 0.3377/0.3320 & 0.927 & 0.207 \\
R1g (true mask)      & 0.8343 (0.7269--0.9270) & 0.9933 & 0.6716 & 0.2071 & 0.3183/0.3145 & 0.927 & 0.075 \\
R2g (mask+ring)      & 0.8507 (0.7500--0.9448) & 0.9940 & 0.5000 & 0.2828 & 0.4175/0.4173 & 0.985 & 0.052 \\
\midrule
\multicolumn{8}{l}{Paired ($\alpha{=}0.01$; $n_{\mathrm{boot}}{=}4952$ for AUROC after discarding single-class}\\
\multicolumn{8}{l}{resamples, 5000 for sens): T1c +0.1315 ($-$0.1208, 0.4187; $p{=}0.324$) underpowered-null;}\\
\multicolumn{8}{l}{T2c $-$0.1940 ($-$0.4193, 0.0651; $p{=}0.117$) sign-flipped vs in-domain T2, CI crosses 0;}\\
\multicolumn{8}{l}{T3c +0.1230 (0.0000, 0.2896; $p{=}0.049$) nominal only; T4c $-$0.3170 ($-$0.4762, $-$0.1423;}\\
\multicolumn{8}{l}{$p{=}0.0004$) \textbf{YES}; T5c $-$0.1570 ($-$0.2537, $-$0.0597; $p{=}0.0028$) \textbf{YES}.}
\end{tabular}
\end{table}

\noindent\emph{Table S10, referenced from main text, `External classification battery' (TNUI-2021 val and the pre-declared pool; families Td and P).}

\begin{table}[!ht]
\centering\scriptsize
\setlength{\tabcolsep}{3pt}
\caption{Phase-3d: TNUI-val ($n{=}274$; 258 malignant/16 benign) and the
pre-declared pooled TNUI test+val ($n{=}413$; 392/21). Locked thresholds from
TN5000-val (raw-raw). Val-panel ECE-pre range 0.24--0.42; pooled panel reports
post-temperature ECE. Replication gate vs locked 3c: zero drift.}
\label{tab:phase3d}
\begin{tabular}{lccccc}
\toprule
Model & AUROC (95\% CI) & AUPRC & sens@spec90 & Brier & ECE post \\
\midrule
\multicolumn{6}{l}{\emph{TNUI-2021 val}}\\
ERM  & 0.6237 (0.4765--0.7670) & 0.9617 & 0.5504 & 0.1879 & 0.2497 \\
R0   & 0.5901 (0.3824--0.7877) & 0.9338 & 0.4031 & 0.3019 & 0.4162 \\
R1p  & 0.5775 (0.3895--0.7484) & 0.9422 & 0.6512 & 0.2434 & 0.3278 \\
R2p  & 0.6008 (0.4193--0.7649) & 0.9520 & 0.6318 & 0.2196 & 0.3121 \\
R1g (true mask)  & 0.6126 (0.4516--0.7640) & 0.9557 & 0.6434 & 0.2475 & 0.3112 \\
R2g (mask+ring)  & 0.5770 (0.4111--0.7370) & 0.9536 & 0.5271 & 0.3059 & 0.3902 \\
\midrule
\multicolumn{6}{l}{\emph{Pooled TNUI test\,+\,val (single pre-declared pool)}}\\
ERM  & 0.5969 (0.4756--0.7106) & 0.9660 & 0.5459 & 0.1907 & 0.2557 \\
R0   & 0.5850 (0.4284--0.7370) & 0.9505 & 0.4056 & 0.3038 & 0.4172 \\
R1p  & 0.6079 (0.4726--0.7385) & 0.9564 & 0.6658 & 0.2300 & 0.3156 \\
R2p  & 0.6144 (0.4725--0.7519) & 0.9615 & 0.6378 & 0.2183 & 0.3188 \\
R1g (true mask)  & 0.6635 (0.5414--0.7780) & 0.9677 & 0.6531 & 0.2339 & 0.3111 \\
R2g (mask+ring)  & 0.6340 (0.5031--0.7591) & 0.9661 & 0.5179 & 0.2981 & 0.3953 \\
\midrule
\multicolumn{6}{l}{Pooled paired family ($\alpha{=}0.01$, 5000 resamples): P1 R1p$-$R0 +0.0227 ($-$0.0718,}\\
\multicolumn{6}{l}{0.1234; $p{=}0.662$) no; P2 ERM$-$R1p $-$0.0118 ($-$0.1264, 0.0991; $p{=}0.852$) no;}\\
\multicolumn{6}{l}{P3 R1g$-$R1p +0.0554 ($-$0.0077, 0.1293; $p{=}0.096$) no; P4 ERM$-$R1g $-$0.0672}\\
\multicolumn{6}{l}{($-$0.1822, 0.0382; $p{=}0.224$) no; P5 ERM$-$R1p sens $-$0.1201 ($-$0.1811, $-$0.0612;}\\
\multicolumn{6}{l}{$p{<}0.0004$) \textbf{YES} --- the 3c operating-point inversion replicates across cohorts.}
\end{tabular}
\end{table}

\noindent\emph{Table S11, referenced from main text, `External classification battery' (DDTI and TNUI-2021 train; families F and E).}

\begin{table}[!ht]
\centering\footnotesize
\caption{Phase-3e arms (locked): DDTI XML-labelled subset ($n{=}189$;
145/44) and TNUI-train ($n{=}966$; 915/51). Conventions as in
Table~\ref{tab:phase3d}. DDTI is at chance for every model, with true-mask
oracle rows \emph{below} 0.5 --- burned-in acquisition annotations are a
documented confound.}
\label{tab:phase3e}
\begin{tabular}{lccccc}
\toprule
Model & AUROC (95\% CI) & AUPRC & sens@spec90 & Brier & ECE post \\
\midrule
\multicolumn{6}{l}{\emph{DDTI (independent institution)}}\\
ERM  & 0.4746 (0.3729--0.5740) & 0.7675 & 0.4828 & 0.2704 & 0.2699 \\
R0   & 0.5116 (0.4196--0.6034) & 0.7923 & 0.1724 & 0.4174 & 0.4364 \\
R1p  & 0.4967 (0.4059--0.5844) & 0.7874 & 0.3379 & 0.3874 & 0.3864 \\
R2p  & 0.5017 (0.4057--0.5909) & 0.7928 & 0.3034 & 0.3696 & 0.3559 \\
R1g (true mask)  & 0.4216 (0.3315--0.5066) & 0.7668 & 0.4276 & 0.3940 & 0.3927 \\
R2g (mask+ring)  & 0.4143 (0.3210--0.5020) & 0.7435 & 0.2759 & 0.4334 & 0.4263 \\
\midrule
\multicolumn{6}{l}{\emph{TNUI-train (heterogeneity probe)}}\\
ERM  & 0.5936 (0.5140--0.6706) & 0.9604 & 0.5519 & 0.1987 & 0.2593 \\
R0   & 0.5472 (0.4556--0.6407) & 0.9505 & 0.4230 & 0.3104 & 0.4108 \\
R1p  & 0.6316 (0.5491--0.7160) & 0.9633 & 0.6776 & 0.2151 & 0.2968 \\
R2p  & 0.6441 (0.5714--0.7170) & 0.9698 & 0.6404 & 0.2136 & 0.2978 \\
R1g (true mask)  & 0.6459 (0.5603--0.7301) & 0.9643 & 0.6656 & 0.2248 & 0.3036 \\
R2g (mask+ring)  & 0.6078 (0.5222--0.6953) & 0.9605 & 0.5148 & 0.3011 & 0.3940 \\
\midrule
\multicolumn{6}{l}{E family (DDTI, $\alpha{=}0.01$, 5000 resamples): E1 R1p$-$R0 $-$0.0162 ($-$0.1064,}\\
\multicolumn{6}{l}{0.0768; $p{=}0.724$); E2 ERM$-$R1p $-$0.0215 ($-$0.1220, 0.0823; $p{=}0.676$); E3 R1g$-$R1p}\\
\multicolumn{6}{l}{$-$0.0749 ($-$0.1498, $-$0.0006; $p{=}0.048$) nominal only --- oracle \emph{worse} than}\\
\multicolumn{6}{l}{predicted ROI on DDTI; E4 ERM$-$R1g +0.0534 ($-$0.0407, 0.1473; $p{=}0.273$);}\\
\multicolumn{6}{l}{\textbf{E5 ERM$-$R1p sens +0.1456 (0.0621, 0.2345; $p{=}0.0012$) YES} --- operating-point}\\
\multicolumn{6}{l}{inversion on the independent cohort.}\\
\multicolumn{6}{l}{F family (TNUI-train): F1 R1p$-$R0 +0.0846 (0.0083, 0.1642; $p{=}0.025$) nominal-}\\
\multicolumn{6}{l}{positive; F2 $-$0.0381 ($p{=}0.326$); F3 +0.0144 ($p{=}0.566$); F4 $-$0.0525 ($p{=}0.222$);}\\
\multicolumn{6}{l}{\textbf{F5 ERM$-$R1p sens $-$0.1259 ($-$0.1628, $-$0.0885; $p{<}0.0004$) YES} --- the grounded}\\
\multicolumn{6}{l}{operating-point advantage replicates in the third TNUI cohort.}
\end{tabular}
\end{table}

\noindent\emph{Figure S1, referenced from main text, `External
classification battery' (external verdict paragraph).}

% FIG 3 float env [H] -- generated from figure_captions.tex; do not hand-edit
\begin{figure}[H]
\centering
\includegraphics{figures/fig3_operating_points.pdf}
\caption{\textbf{The operating point is the unit of transfer.}
Sensitivity at each model's frozen specificity-$\geq$0.90 threshold
across cohorts (image level except ThyUS2Path, which aggregates to case
level under the audit-gated protocol). Fine-tuned ERM leads in-domain
but collapses under shift (0.802 to 0.013 on ThyUS2Path); the grounded
frozen head (R1p) degrades least on every external cohort except the
documented DDTI counterexample and retains the lead at case level
(0.159 vs 0.013; McNemar 77 vs 1 $p{<}0.0004$).}
\label{fig:operating_points}
\end{figure}



\noindent\emph{Table S12, referenced from main text, `Joint partial-supervision fine-tuning' (in-domain levels).}

\begin{table}[!ht]
\centering\scriptsize
\setlength{\tabcolsep}{3pt}
\caption{Phase-4 in-domain TN5000 test ($n{=}950$; locked, protocol v2).
Columns as Table~\ref{tab:phase3}; JOINT-row CIs from this phase's
2,000-replicate case bootstrap. Both segmentation non-collapse gates passed
(TN3K-test Dice 0.8258 JOINT-g / 0.8165 JOINT-img against the
$\geq 0.7966$ gate). Context rows are the locked Phase-3 values
(Table~\ref{tab:phase3}).}
\label{tab:phase4_indomain}
\begin{tabular}{lcccccc}
\toprule
Model & AUROC (95\% CI) & AUPRC & sens@spec90 & Brier & ECE pre/post & $T_{\mathrm{val}}$ \\
\midrule
JOINT-g   & 0.8698 (0.8426--0.8939) & 0.9423 & 0.5642 & 0.1573 & 0.1971/0.1682 & 0.523 \\
JOINT-img & 0.7207 (0.6824--0.7591) & 0.8460 & 0.1962 & 0.2103 & 0.2241/0.1927 & 0.200 \\
\midrule
\multicolumn{7}{l}{\emph{Locked Phase-3 context rows}}\\
ERM whole-image (ResNet34 ft) & 0.9401 (0.9232--0.9553) & 0.9738 & 0.8023 & 0.0878 & 0.0893/0.0634 & 1.969 \\
R0 image | linear             & 0.8534 (0.8271--0.8798) & 0.9358 & 0.6479 & 0.1520 & 0.1278/0.1244 & 0.900 \\
R1p lesion | linear           & 0.8902 (0.8670--0.9122) & 0.9519 & 0.7273 & 0.1515 & 0.1788/0.1734 & 0.847 \\
R2p lesion+ring | linear      & 0.8829 (0.8584--0.9061) & 0.9464 & 0.6869 & 0.1538 & 0.1721/0.1705 & 0.956 \\
R1g GT box | linear (oracle)  & 0.9105 (0.8899--0.9299) & 0.9635 & 0.8254 & 0.1260 & 0.1339/0.1302 & 0.927 \\
R2g GT box+ring | linear (oracle) & 0.9029 (0.8807--0.9239) & 0.9569 & 0.7763 & 0.1449 & 0.1750/0.1721 & 0.927 \\
\bottomrule
\end{tabular}
\end{table}

\noindent\emph{Table S13, referenced from main text, `Joint partial-supervision fine-tuning' (external levels).}

\begin{table}[!ht]
\centering\scriptsize
\setlength{\tabcolsep}{3pt}
\caption{Phase-4 external battery (locked; thresholds from each model's own
TN5000-val raw probabilities, raw-raw convention). Segmentation QC for the
JOINT-g trunk: external Dice 0.7201 (TNUI-test), 0.7651 (TNUI-val), 0.7543
(TNUI-train), 0.6738 (DDTI; TN3K gate 0.8258); empty-mask rate $\leq 0.3\%$
on every cohort. Frozen rows (ERM, R1p) replicate the locked Phase-3c/3d/3e
external values exactly. $^\dagger$descriptive cohorts under the locked
pooling policy (inference only within the pooled family, family H,
the main-text Phase-4 family-G/H table).}
\label{tab:phase4_external}
\begin{tabular}{llcccccccc}
\toprule
 & & \multicolumn{4}{c}{AUROC} & \multicolumn{4}{c}{sens@spec90} \\
\cmidrule(lr){3-6}\cmidrule(lr){7-10}
Cohort & $n$ (mal) & ERM & JOINT-g & JOINT-img & R1p & ERM & JOINT-g & JOINT-img & R1p \\
\midrule
TNUI-test  & 139 (134) & 0.5172 & 0.7075 & 0.6493 & 0.7119 & 0.5373 & 0.6119 & 0.1343 & 0.6940 \\
TNUI-val   & 274 (258) & 0.6237 & 0.5523 & 0.5337 & 0.5775 & 0.5504 & 0.6124 & 0.1434 & 0.6512 \\
TNUI-train$^\dagger$ & 966 (915) & 0.5936 & 0.5845 & 0.5270 & 0.6316 & 0.5519 & 0.6251 & 0.1574 & 0.6776 \\
DDTI$^\dagger$       & 189 (145) & 0.4746 & 0.5210 & 0.5296 & 0.4967 & 0.4828 & 0.2621 & 0.0414 & 0.3379 \\
\midrule
\multicolumn{10}{l}{Derived pooled (test+val) operating points: sens@spec90
R1p $\approx 0.666$ $>$ JOINT-g $\approx 0.612$ $>$ ERM $\approx 0.546$.}\\
\bottomrule
\end{tabular}
\end{table}

\noindent\emph{Table S14, referenced from main text, `Joint partial-supervision fine-tuning' (selective-prediction rows).}

\begin{table}[!ht]
\centering\small
\caption{Phase-4 selective prediction (Gap 6; locked). Coverage = the
most-confident-retained fraction at the frozen raw validation threshold
(raw-raw; the validation-fitted temperature, $T{=}0.523$, supplies the
confidence margin only --- no threshold or temperature is re-estimated).
\textit{mal.\ FNR} = malignant false-negative rate among retained cases
(equal to $1-$sens@spec90 row-wise); \textit{mal.\ share} = fraction of the
cohort's malignant cases falling inside the retained confident set.
Companion figure: \texttt{phase4\_selective.png} (both curves monotone in
coverage).}
\label{tab:phase4_selective}
\begin{tabular}{lcccc}
\toprule
Coverage & $n$ retained & mal.\ FNR & sens@spec90 & mal.\ share retained \\
\midrule
\multicolumn{5}{l}{\emph{TN5000-test (in-domain, $n{=}950$, 693 malignant)}}\\
100\% & 950 & 0.4358 & 0.5642 & 1.000 \\
90\%  & 855 & 0.3673 & 0.6327 & 0.8918 \\
80\%  & 760 & 0.2826 & 0.7174 & 0.7864 \\
70\%  & 665 & 0.1602 & 0.8398 & 0.6667 \\
60\%  & 570 & 0.1490 & 0.8510 & 0.5714 \\
50\%  & 475 & 0.1262 & 0.8738 & 0.4690 \\
\midrule
\multicolumn{5}{l}{\emph{TNUI pooled (test+val; external, $n{=}413$, 392 malignant)}}\\
100\% & 413 & 0.3878 & 0.6122 & 1.000 \\
90\%  & 372 & 0.3220 & 0.6780 & 0.9031 \\
80\%  & 330 & 0.2357 & 0.7643 & 0.8010 \\
70\%  & 289 & 0.1594 & 0.8406 & 0.7041 \\
60\%  & 248 & 0.1398 & 0.8602 & 0.6020 \\
50\%  & 206 & 0.1026 & 0.8974 & 0.4974 \\
\bottomrule
\end{tabular}
\end{table}

\noindent\emph{Figure S2, referenced from main text, `Joint
partial-supervision fine-tuning' (selective-prediction paragraph).}

% FIG 4 float env [H] -- generated from figure_captions.tex; do not hand-edit
\begin{figure}[H]
\centering
\includegraphics{figures/fig4_selective_coverage.pdf}
\caption{\textbf{Selective prediction is referral-shaped, in-domain and
externally.} (a) Sensitivity at the frozen operating point as
lower-confidence cases are referred (coverage = most-confident retained
fraction; confidence from the validation-fitted temperature margin, no
refitting). (b) Fraction of each cohort's malignant cases inside the
retained set. At 50\% coverage, sensitivity reaches 0.874 (TN5000-test)
and 0.897 (pooled TNUI), and 86--87\% of the frozen point's malignant
misses fall into the referred half.}
\label{fig:selective}
\end{figure}



\noindent\emph{Table S15, referenced from main text, `Case-level external validation on ThyUS2Path' (per-model case metrics).}

\begin{table}[!ht]
\centering\scriptsize
\setlength{\tabcolsep}{3pt}
\caption{ThyUS2Path case-level external battery (locked; $n{=}839$ cases,
522 malignant; unit = case, score = mean of per-image calibrated
probabilities; thresholds from each model's own TN5000-val raw
probabilities, raw-raw convention). CIs: 2{,}000-replicate case
bootstrap. Oracle regimes (R1g/R2g) not computable --- ThyUS2Path
provides no lesion annotation (pre-declared exclusion). Companion file:
\texttt{phase5\_main\_table.csv}.}
\label{tab:phase5_main}
\begin{tabular}{lcccccc}
\toprule
Model & AUROC (95\% CI) & AUPRC & sens@spec90 & Brier & ECE pre/post & $T_{\mathrm{val}}$ \\
\midrule
ERM whole-image (ResNet34 ft) & 0.5479 (0.5082--0.5868) & 0.6662 & 0.0134 & 0.2873 & 0.2813/0.2145 & 1.969 \\
R0 image | linear             & 0.5400 (0.4992--0.5791) & 0.6408 & 0.0460 & 0.3486 & 0.3039/0.3139 & 0.927 \\
R1p lesion | linear           & 0.5081 (0.4686--0.5494) & 0.6300 & 0.1590 & 0.3423 & 0.2761/0.2942 & 0.847 \\
R2p lesion+ring | linear      & 0.5211 (0.4805--0.5619) & 0.6392 & 0.0881 & 0.3337 & 0.2771/0.2852 & 0.927 \\
JOINT-g                       & 0.5112 (0.4715--0.5506) & 0.6342 & 0.0613 & 0.3200 & 0.2020/0.2607 & 0.523 \\
JOINT-img                     & 0.4780 (0.4387--0.5201) & 0.6078 & 0.0153 & 0.3198 & 0.1506/0.2738 & 0.200 \\
\bottomrule
\end{tabular}
\end{table}

\noindent\emph{Table S16, referenced from main text, `Case-level external validation on ThyUS2Path' (domain-shift matrix).}

\begin{table}[!ht]
\centering\small
\caption{Domain shift (locked): TN5000-test AUROC (source) versus
ThyUS2Path case-level AUROC (external). Source values are re-derived in
the Phase-5 environment, not copied: ERM re-hits the locked 0.9401 sanity
gate; R1p reproduces 0.8889 (within the 0.002 integrity tolerance of the
locked 0.8902 reference); the JOINT rows match
Table~\ref{tab:phase4_indomain} exactly.}
\label{tab:phase5_shift}
\begin{tabular}{lccc}
\toprule
Model & AUROC TN5000-test & AUROC ThyUS2Path & Shift drop \\
\midrule
ERM       & 0.9401 & 0.5479 & 0.3921 \\
R0        & 0.8524 & 0.5400 & 0.3124 \\
R1p       & 0.8889 & 0.5081 & 0.3808 \\
R2p       & 0.8841 & 0.5211 & 0.3630 \\
JOINT-g   & 0.8698 & 0.5112 & 0.3586 \\
JOINT-img & 0.7207 & 0.4780 & 0.2427 \\
\midrule
\multicolumn{4}{l}{Derived mean shift drop: 0.3416 AUROC (range 0.2427--0.3921).}\\
\bottomrule
\end{tabular}
\end{table}

\noindent\emph{Table S17, referenced from main text, `Case-level external validation on ThyUS2Path' (first-listed-image view).}

\begin{table}[!ht]
\centering\small
\caption{First-listed-image sensitivity view (descriptive; locked). Each
case is scored by its first-listed image only --- identical frozen
thresholds, no case aggregation --- demonstrating that the collapse of
Table~\ref{tab:phase5_main} is not an artefact of the case-mean unit.}
\label{tab:phase5_firstimage}
\begin{tabular}{lcc}
\toprule
Model & AUROC (first image) & sens@spec90 (first image) \\
\midrule
ERM       & 0.5328 & 0.1648 \\
R0        & 0.5006 & 0.1475 \\
R1p       & 0.4877 & 0.2375 \\
R2p       & 0.4965 & 0.2069 \\
JOINT-g   & 0.4813 & 0.1743 \\
JOINT-img & 0.4974 & 0.0632 \\
\bottomrule
\end{tabular}
\end{table}

\noindent\emph{Table S18, referenced from main text, `Lesion-grounded explainability battery' (localisation ratios, in-box mass, entropy, pointing games).}

\begin{table}[!ht]
\centering\small
\caption{Phase-6 plausibility against TN5000 expert boxes ($n{=}950$;
all images boxed; localisation ratio $=$ in-box share $/$ box-area
share; in-box mass $=$ share of the $\ell_1$-normalised explanation
inside the box; hit $=$ argmax inside box (pointing game, descriptive);
entropy normalised to the 384$^2$ canvas. Score-CAM computed on the
pre-registered seeded $n{=}150$ subsample (top-64 channels by
$|\alpha|$). All quantities computed once, seed 42, no refit.}
\label{tab:phase6_plausibility}
\begin{tabular}{lcccc}
\toprule
Explanation & loc.\ ratio mean (median) & in-box mass & pointing hit & norm.\ entropy \\
\midrule
R1p-ELA (grounded, exact)      & 41.63 (27.83) & 0.796 & 0.637 & 0.717 \\
ERM Grad-CAM (layer4)          & 3.60 (4.44)   & 0.146 & 0.514 & 0.928 \\
ERM Eigen-CAM                  & 4.68 (4.80)   & ---   & 0.736 & ---   \\
ERM Score-CAM ($n{=}150$)      & 3.57 (3.57)   & ---   & 0.580 & ---   \\
\bottomrule
\end{tabular}
\end{table}

\noindent\emph{Table S19, referenced from main text, `Lesion-grounded explainability battery' (deletion/insertion areas incl. random-control mechanism).}

\begin{table}[!ht]
\centering\small
\caption{Phase-6 faithfulness on TN5000-test (deletion/insertion under
per-image blur $\sigma{=}8$, 8 steps, raw-raw probabilities, masks fixed
post-fallback; gap $=$ trapezoid-area difference against the per-image
seeded random-ordering control; insertion descriptive on the seeded
$n{=}300$ subset). ThyUS-600 per-model curves are archived in
\texttt{phase6\_panel\_thyus.npz}; the inferential external statement
(K3) is in the main-text family-K table.}
\label{tab:phase6_faithfulness}
\begin{tabular}{lcccc}
\toprule
Model & del.\ AUC guided & del.\ AUC random & deletion gap & insertion gap ($n{=}300$) \\
\midrule
R1p       & 0.3334 & 0.4986 & 0.1652 & 0.0351 \\
ERM       & 0.3315 & 0.5564 & 0.2249 & 0.2043 \\
JOINT-g   & 0.3886 & 0.4908 & 0.1022 & 0.0356 \\
JOINT-img & 0.4405 & 0.4815 & 0.0410 & 0.0144 \\
\bottomrule
\end{tabular}
\end{table}

\noindent\emph{Table S20, referenced from main text, `Lesion-grounded explainability battery' (noise/affine stability rows).}

\begin{table}[!ht]
\centering\small
\caption{Phase-6 stability under perturbation (mean / median Spearman
against the unperturbed explanation; seeded 160-image probes per cell;
pre-declared descriptive, non-gated).}
\label{tab:phase6_stability}
\begin{tabular}{llccc}
\toprule
Cohort & Model & Perturbation & mean $\rho$ & median $\rho$ \\
\midrule
TN5000-test & R1p & noise $\sigma{=}0.02$  & 0.9208 & 0.9972 \\
TN5000-test & R1p & affine ($\pm 3^{\circ}$, $\pm 5$px) & 0.7713 & 0.8815 \\
TN5000-test & ERM & noise $\sigma{=}0.02$  & 0.8317 & 0.8865 \\
TN5000-test & ERM & affine ($\pm 3^{\circ}$, $\pm 5$px) & 0.8709 & 0.9245 \\
ThyUS-600   & R1p & noise $\sigma{=}0.02$  & 0.7990 & 0.9782 \\
ThyUS-600   & R1p & affine ($\pm 3^{\circ}$, $\pm 5$px) & 0.6402 & 0.7909 \\
ThyUS-600   & ERM & noise $\sigma{=}0.02$  & 0.7545 & 0.7879 \\
ThyUS-600   & ERM & affine ($\pm 3^{\circ}$, $\pm 5$px) & 0.8103 & 0.8574 \\
\bottomrule
\end{tabular}
\end{table}

\noindent\emph{Table S21, referenced from main text, `Lesion-grounded explainability battery' (randomisation cascade stages).}

\begin{table}[!ht]
\centering\small
\caption{Phase-6 randomisation cascade (mean Spearman between the
original and randomised explanations over the seeded 160-image TN5000
probe). Pre-registered self-gate $|\rho|{<}0.2$ at each model's final
available stage: ERM \textbf{pass}, R1p \textbf{fail} --- reported as
failed. The fc-only row shows shallow readout randomisation cannot reach
the bar for either model; follow-up diagnostics (decision rules
committed pre-run): resampled-vs-resampled floor 0.0920 (elevation of
the retained correlation $+0.2072$ above floor), per-feature attribution
decorrelation complete (Spearman $-0.0150$). Originals verified
untouched by checksum.}
\label{tab:phase6_cascade}
\begin{tabular}{lccc}
\toprule
Stage & mean $|\rho|$ & mean $\rho$ & $n$ \\
\midrule
ERM fc                          & 0.2948 & $-0.0618$ & 160 \\
ERM fc+layer4                   & 0.1492 & $+0.1287$ & 160 \\
ERM fc+layer4+layer3 (final)    & 0.1205 & $-0.0135$ & 160 \\
R1p head weights (final)        & 0.2991 & $+0.2982$ & 160 \\
\bottomrule
\end{tabular}
\end{table}

\noindent\emph{Table S22: family-K secondary/tertiary contrasts,
relocated row-for-row from the main-text family-K table (referenced
from main text, `Lesion-grounded explainability battery'). Conventions
as in the main-text table (image-level paired bootstrap, 5{,}000
resamples, seed 42; Bonferroni $\alpha{=}0.0125$; $p{=}0.0$ typeset
$p{<}0.0004$).}

\begin{table}[!ht]
\centering\scriptsize
\setlength{\tabcolsep}{3pt}
\caption{Supplementary family-K secondary contrasts (relocated rows;
locked values). ``YES (rev.)'' = significant in the
pre-registered-opposite direction (see main text).}
\label{tab:sm_familyk_sec}
\begin{tabular}{llllccl}
\toprule
Test & Metric & A & B & $\Delta$(A$-$B) (95\% CI) & Sig. & Note \\
\midrule
K2 & deletion gap       & R1p     & ERM         & $-0.0598$ ($-0.0761..-0.0437$)  & YES (rev.) & faithfulness, TN5000 \\
K3 & deletion gap       & R1p     & ERM         & $-0.0339$ ($-0.0511..-0.0171$)  & YES (rev.) & faithfulness, ThyUS-600 \\
K4 & deletion gap       & JOINT-g & JOINT-img   & $+0.0613$ (0.0560..0.0667)      & YES & grounding, G3 mirror \\
\bottomrule
\end{tabular}
\end{table}

\section{Narrative extracts preserved from the pre-compression main text}
\label{sm:extracts}

\emph{Sentences below are reproduced verbatim from the pre-compression
main text wherever a reported detail has no other Supplementary home,
so the compression removes nothing from the submission package.}
\par
\paragraph{From Material and Methods (architecture and preprocessing), verbatim.}
The working resolution is $384\times384$; the prototype module uses levels
$\{e_2, e_3, e_4\}$ (strides 8/16/32; 128/256/512 channels) and
descriptors pool $\{e_3, e_4\}$. Images are letterboxed; masks share the
geometric transform (nearest-neighbour) and are binarised at grey level
$\geq 127$, selected by visual audit against $>0$ and $\geq 250$
alternatives; inputs are ImageNet-normalised. Score maps are bilinearly
resized to the stride-8 grid ($48\times48$) and stacked as $2L$ channels
(fragment, equation region).

\paragraph{From Material and Methods (episode and budget mechanics), verbatim.}
Meta-training samples 300 episodes per epoch from the source pool,
re-drawing $K_m$ uniformly from $\{1,5,10,20\}$ at every episode;
validation uses 100 episodes and evaluation 200 per condition.
Descriptors are L2-normalised, and per-class means over nested per-class
draws from the training split ($K_c \in
\{5,10,20,50,\allowbreak 200,\mathrm{all}\}$, 25 resamples per finite
budget, fixed seed) are scored by the cosine margin, reported as mean
$\pm$ standard deviation of test AUROC across draws; in-domain only.

\paragraph{From Material and Methods (decision conventions), verbatim.}
Four conventions were fixed after the paired-analysis audit: (1)
thresholds and temperatures come from the source validation split only,
frozen before any test or external evaluation; (2) the raw-raw
convention never mixes raw-threshold binary decisions with calibrated
probabilities (rank metrics are convention-invariant); (3) confidence
intervals are case-level bootstrap, 2,000 resamples, percentile
2.5/97.5; (4) externally, calibration metrics are descriptive only ---
any deployment interpretation requires local recalibration, stated
rather than performed.

\paragraph{From Material and Methods / Introduction (dataset roster), composed verbatim.}
TN-SCUI2020 provides class labels with delineations (dataset record in
the main-text reference list); TN-SCUI2020 remains an optional extension
pending authorised access.

\paragraph{From Results (Phase 1), verbatim.}
TTA effects are small and sign-mixed --- the largest shift is $+1.11$
Dice points on TNUI-all, and DDTI moves slightly in the opposite
direction --- and DDTI is visibly the harder masked cohort (sensitivity
0.83, median HD95 roughly doubled), consistent with its lower zero-shot
prototype scores. Best validation episodic Dice 0.7776 (epoch 13; early
stop 21; 37.7 min).
DDTI is the hardest masked cohort for both paradigms (supervised 0.7396;
zero-shot 0.6782), consistent with its doubled boundary error (median
HD95 ${\approx}22$ vs 10--11 px) and burned-in annotations.

\paragraph{From Results (Phase 3b, error propagation), verbatim.}
Error propagation is diffuse rather than confidence-localised: the
oracle-minus-predicted AUROC gap is modest in the lowest-confidence
quartiles, vanishes in Q3, and is descriptively largest in the
\emph{highest}-confidence quartile ($+0.0589$), the mean gap
($\approx$0.022) closely matching the global T3 contrast. Segmentation
error therefore imposes a diffuse $\approx$2-point tax on grounded
features, and there is no confidence-gated ``safe subgroup'' in which
predicted masks are output-equivalent to oracle boxes.

\paragraph{From Results (Phase 3, label-budget sweep), verbatim.}
The class-prototype label-budget sweep ($K$ per class
$\in \{5,10,20,50,200,\mathrm{all}\}$, 25 draws) carries the
classification-side annotation-efficiency claim: grounding dominates the
identical whole-image budget at every $K$ ($+0.046$ to $+0.096$ AUROC);
R1p saturates by $K{=}50$ (0.6383 versus 0.6379 full-budget) and stays
within 0.013--0.029 of the oracle, while R0 has not saturated even at
$K{=}200$; and the ring is null-or-negative throughout (R2p $\leq$ R1p,
R2g $\leq$ R1g at every $K$), corroborating the T4 null under a
different head family. Absolute levels sit far below the linear heads
(0.59--0.66 vs 0.85--0.91), as expected for a cosine prototype rule on
frozen descriptors; the linear head remains the pre-declared primary
classifier.

\paragraph{From Results (Phase 4, external family and abstention), verbatim.}
The primary external test is nominal only: on the single locked TNUI
test+val pool JOINT-g exceeds ERM's sensitivity at the frozen operating
point by 6.6 points (H1: $+0.0661$, 95\% CI 0.005..0.128, $p{=}0.033$,
beyond the family bar; McNemar 86 vs 60). The rank test is a flat null
(H3: $-0.011$, $p{=}0.858$), and H2 is significant: joint grounding
transfers its operating point \emph{worse} than the frozen R1p by 5.3
points of pooled sensitivity (H2: $-0.0534$, CI $-0.097..-0.010$,
$p{=}0.011$). The joint segmenter travels adequately (external Dice
0.674--0.765, empty-mask $\leq 0.3\%$), so the deficit lives in
classification head-room, not mask collapse. JOINT-g had not saturated
at the pre-registered 15-epoch cap (validation AUROC was still
improving; the cap was not extended), so family-G magnitudes are
conservative. Retaining the most-confident 70\% of TN5000-test cases
lifts JOINT-g's sensitivity at spec90 from 0.564 to 0.840, and at 50\%
coverage the curve reaches 0.874 in-domain and 0.897 on the external
TNUI pool. The effect is referral-shaped, not accuracy-shaped:
approximately 86--87\% of the frozen point's malignant misses sit inside
the abstained lower-confidence half on both panels (261 of 302 misses
in-domain; 132 of 152 on pooled TNUI). No threshold or temperature was
re-estimated for these numbers (raw-raw; the validation-fitted
$T{=}0.523$ supplies only the confidence margin).

\paragraph{From Results (Phase 5, audit and cohort assembly), verbatim.}
Label coverage was complete for Batch 1 (6{,}005/6{,}005) and
near-complete for Batch 2 (2{,}495/2{,}503), giving an overall parse
rate of 0.9991 (8{,}500/8{,}508); the surviving join matched the
corpus's own image map exactly (5{,}893/5{,}893 Batch-1 stems). Assembly
gave 842 cases (601 + 241), and conservative whole-case near-duplicate
screening excluded three, leaving 839 cases and 8{,}465 images
(598 + 241), of which 522 are malignant (prevalence 0.622, inside the
$[0.20, 0.85]$ gate). Frozen-row integrity held exactly (ERM re-hit
0.9401; the retrained R1p head reproduced 0.8889 against the locked
0.8902, within the 0.002 integrity tolerance; both JOINT arms matched to
$\leq 10^{-6}$, temperatures 0.523/0.200), and segmentation QC was
healthy (empty masks on 10.9\% of images, 926/8{,}465; JOINT-g decoder
3.5\% (293/8{,}465)). Mean predicted nodule areas were 4.1\% and 5.9\%
of the frame respectively.

\paragraph{From Results (Phase 5, collapse and calibration detail), verbatim.}
Only ERM's interval clears 0.5 --- by eight thousandths. AUPRC spans
0.608--0.666 against the 0.622 prevalence, and shift drops from source
are 0.31--0.39 AUROC for every linear head and 0.39 for ERM (drop
0.3921, the largest); even JOINT-img, modest in-domain (0.7207), loses
0.24, and the first-list-image view shows the same pattern
(0.481--0.533). Calibration does not travel either: source-fitted
temperatures \emph{worsen} external ECE for five of six models
(post-calibration ECE 0.26--0.31; only ERM improves, 0.2813 $\to$
0.2145), and Brier scores (0.29--0.35) sit far above the source regime.
In integers: of the 83 malignant cases R1p catches at the frozen point,
ERM misses 77; of the 7 ERM catches, R1p misses 1 --- an eleven-fold
relative advantage, and the external case-level replication of the
internal operating-point result.

\paragraph{From Results (Phase 6, integrity detail), verbatim.}
The seeded ThyUS-600 first-image panel ($n{=}600$, 376 malignant) was
drawn from the licensed 839-case cohort and triangulated: 522 malignant,
8{,}465 images, prevalence 0.6222. Live probabilities reproduce the
Phase-3/4 archives (ERM 0.9404 vs 0.9401; R1p 0.8889 vs 0.8889; JOINT
arms to ${\leq}2.25{\times}10^{-5}$ max $|\Delta p|$ against the
$10^{-4}$ gate), and the 32-image ELA exactness probe confirmed the
closed-form logit decomposition (max deviation R0 $1.6{\times}10^{-6}$,
R1p $2.9{\times}10^{-7}$, R2p $2.7{\times}10^{-7}$, JOINT-g
$7.8{\times}10^{-8}$, JOINT-img $2.6{\times}10^{-8}$; gate $10^{-3}$).

\paragraph{From Results (Phase 6, descriptives and faithfulness detail), verbatim.}
R1p-ELA attains a localisation ratio of 41.6 (median 27.8) where
Grad-CAM attains 3.6. The decomposition marks R1p's precision as
segmenter-grade (79.6\% of attribution mass inside the box, normalised
entropy 0.717) against ERM-family diffusion (in-box mass 0.146, entropy
0.928). The pointing game cuts the other way descriptively (Eigen-CAM
0.736 $>$ R1p 0.637 $>$ Score-CAM 0.580 $>$ Grad-CAM 0.514): where a
diffuse map happens to peak inside the box it scores like a concentrated
one; the ratio and in-box-mass statistics instead report where the
\emph{bulk} of the explanation lies. Under shift, ERM's attention
deserts the nodule: Grad-CAM's energy inside the frozen predicted mask
falls from 0.1253 in-domain (${\approx}3.2\times$ the mask's area share)
to 0.0549 on ThyUS-600 (${\approx}1.3\times$). The reversal is carried
by the random-order control, not by the guided curves:
explanation-guided deletion degrades both models essentially identically
(areas 0.3334 vs 0.3315), while random deletion erodes the diffuse
image-wide classifier more slowly (random-control areas 0.4986 vs
0.5564). Insertion (descriptive, $n{=}300$) favours ERM (0.2043 vs
0.0351) with the same proviso, compounded by the out-of-distribution
nature of sharp-lesion-on-blur hybrids for a mask-conditioned head.

\paragraph{From Results (Phase 6, joint ladder, stability, cascade diagnostics), verbatim.}
The deletion-gap ladder rises monotonically with the degree of built-in
grounding (JOINT-img 0.041 $<$ JOINT-g 0.102 $<$ R1p 0.165). Under
intensity noise R1p-ELA is the more stable explanation in both domains
(mean Spearman 0.921 vs 0.832 on TN5000; 0.799 vs 0.755 on ThyUS-600),
while affine warps invert the ordering (0.771 vs 0.871; 0.640 vs
0.810) --- the two rows are the same concentration geometry read in
opposite directions. ERM decorrelates below the bar only at its deepest
stage (0.1205), and at the matching shallow stage (fc only, 0.2948) it
fails identically: a 0.2 threshold at readout depth is unattainable for
either model, rather than evidence of explanation defect. Independent
resampled heads correlate at only 0.0920 on the same images, so the
retained 0.2991 sits $+0.2072$ above that floor: the pixel map is
\emph{not} head-agnostic geometry. Meanwhile, the per-feature logit
attribution decorrelates completely under head resampling (Spearman
$-0.0150$, $|{-}|$ 0.0191; Pearson $-0.0016$).

\paragraph{From Discussion (mechanism and referral), verbatim.}
The parameter-free variant collapses to 0.47--0.57 Dice under zero-shot
transfer, so the contribution sits in the decoder, not in
nearest-prototype matching of the PANet kind. A small referral stream
(${\approx}30$\% of cases) removes the majority of the frozen point's
false negatives without touching a single parameter. The descriptives
cohere: 79.6\% of the grounded map's mass lies inside the expert box
versus 14.6\% for Grad-CAM. Committed diagnostics show the retention is
neither head-agnostic geometry (random-vs-random floor 0.092; retained
elevation $+0.2072$ above it) nor a broken explanation (per-feature
attribution decorrelates completely).

\paragraph{From Results (Phase 4, grounding signal), verbatim.}
Under identical multi-task conditions the grounded head beats the
whole-image head by 14.9 points (G3: $+0.1489$, CI 0.1136..0.1848).



\end{document}
